\section{Accessibility without an antibody: the cylinder surrogate}\label{sec:noab}
One of the strictest filters in the DP3 setting concerns antibody accessibility: we place the
real antibody where it sits in the crystal complex and reject any design with scaffold mass in
its path. That filter is only available where a real antibody exists (all of DP3, most of DP4).
Most targets, though, have no solved antibody --- the polyclonal arm of DP4 tiles whole
antigens with no known binder (Section~\ref{sec:intro}) --- so there is nothing to clash
against. We recover the filter with a native-antigen-aware cylinder surrogate: a cylinder
fitted to the epitope plane along the antibody-approach normal that counts non-epitope scaffold
mass intruding into the volume an approaching antibody would need (Fig.~\ref{fig:cylinder}).

\subsection{Definition}\label{sec:cyldef}
With no antibody available, we approximate the volume an antibody would occupy with a
cylinder fitted to the epitope. The epitope C$\alpha$ coordinates $\{x_j\}_{j=1}^{m}$ define
a plane: the centroid is $c=\tfrac{1}{m}\sum_j x_j$, and the plane normal $\hat n$ is the
smallest-variance singular vector of the centered coordinates $\{x_j-c\}$ (the antibody
approach direction), oriented to point away from the antigen body. The cylinder base is
$b = c + d_0\,\hat n$ with offset $d_0=\SI{-4}{\angstrom}$. A scaffold C$\alpha$ at position
$p$ lies inside the cylinder when, writing $v=p-b$, its axial and radial coordinates satisfy
\begin{equation}\label{eq:cyl}
0 \;\le\; v\cdot\hat n \;\le\; H
\qquad\text{and}\qquad
\bigl\lVert\, v-(v\cdot\hat n)\,\hat n \,\bigr\rVert \;\le\; R,
\end{equation}
with height $H=\SI{40}{\angstrom}$ and radius $R=\SI{16}{\angstrom}$. The plain count is the
number of non-epitope scaffold C$\alpha$ satisfying~\eqref{eq:cyl}; the native-aware count
additionally drops any such atom within \SI{1}{\angstrom} of a native-antigen heavy atom,
i.e.\ sitting where the native antigen already sits, which a real antibody tolerates. The
geometry core is \texttt{episcaf\_analysis/native\_cylinder\_core.py} (self-tested).

\begin{figure}[h]\centering
\figorbox{epitope_cylinder.png}{0.5}
\caption{The accessibility cylinder. Epitope residues (red) sit on the scaffold (magenta);
the epitope C$\alpha$ atoms define a plane (orange grid) whose outward normal (blue arrow,
the antibody-approach direction) orients the cylinder (cyan). Scaffold C$\alpha$ atoms
inside the cylinder count against accessibility. Rendered in VMD by
\texttt{known\_antigen/analysis/cylinder\_vis/visualize\_cylinder.tcl}.}
\label{fig:cylinder}
\end{figure}

\subsection{The native-aware carve removes false penalties}\label{sec:carve}
A naive cylinder over-penalizes, because some scaffold mass sits exactly where the native
antigen already sits, a volume a real antibody clearly tolerates since it binds the native
antigen there. The native-aware carve fixes this: among truly clash-free designs, we drop
scaffold C$\alpha$ atoms within \SI{1}{\angstrom} of a native-antigen heavy atom. Of the
\num{3186} DP3 designs with zero true antibody clash, the number the plain cylinder flags as
heavily penalized (count $\ge 10$) falls from \num{1112} to \num{384} once the native
footprint is carved out, so \textbf{\num{728} false-positive penalties are removed}, while
genuinely clashing designs are left essentially untouched (Figs.~\ref{fig:fp}
and~\ref{fig:fpsel}).

\begin{figure}[h]\centering
\figorbox{fp_reduction.png}{0.8}
\caption{Cylinder count among clash-free DP3 designs (zero AlphaFold3 antibody clash,
$n=\num{3186}$), before (grey) and after (red) subtracting the native-antigen volume. The
heavily penalized tail (count $\ge 10$) collapses from \num{1112} to \num{384}. Regenerated
by \texttt{episcaf\_analysis/viz/plot\_fp\_reduction.py --native metrics\_native\_cyl.csv
--mode dist}.}
\label{fig:fp}
\end{figure}

\begin{figure}[h]\centering
\figorbox{fp_selectivity.png}{0.8}
\caption{The carve removes flags selectively. Fraction of designs flagged (cylinder
$\ge 10$) across bins of true antibody clash, before versus after the carve. The gap is
large where there is no real clash (false positives removed) and closes where the clash is
real (true signal retained). \texttt{--mode selectivity}.}
\label{fig:fpsel}
\end{figure}

\paragraph{A pipeline constraint worth recording: crystal gaps.}
We measure 12-mer epitope fidelity by aligning the AlphaFold3 epitope onto the native crystal
epitope, matched by residue number (\texttt{build\_12mer\_metrics.py}). When a tile's epitope
residues are unresolved in the crystal (a gap in the deposited structure), the native
C$\alpha$ anchors are missing, the Kabsch alignment cannot be formed, and the tile is dropped
with status \texttt{crystal\_epi\_incomplete}. So which 12-mer epitopes are scorable at all
depends on crystal coverage: an antigen with gaps over the tiled region contributes fewer (or
no) scorable tiles, regardless of design quality. This is exactly what sets the
epitopes-kept counts in Table~\ref{tab:tilecount} (6M0J and 4WAT lose 15 and 40 tiles), and the
same constraint reappears in the per-island production run (Section~\ref{sec:metricsresults}).

\subsection{Interpreting the cylinder: where passing designs land}\label{sec:cylinterp}
The cylinder is easiest to read by overlaying, for each metric, the full design population
against the designs we would actually advance (Fig.~\ref{fig:overlay}). For DP3 the advanced
set is the four-filter passes (\num{751} designs); for DP4 it is the top three per epitope by
composite score. On every geometry and confidence metric the passing designs concentrate
where we expect: low RMSD, low PAE, high pTM. The cylinder column is the interesting one. The
passing DP3 designs, which we know the antibody accepts, sit at low cylinder counts (median
near 10), while the DP4 top-3 sit higher (medians near 25--30). This is the scaffold-mass
effect again, not a sign the DP4 designs are worse: with a 12-residue epitope on a
$\sim$104-residue design, more non-epitope scaffold mass falls in the fixed cylinder even
among the designs we would pick. This is why we lean on the native-aware carve and per-antigen
scoring rather than one absolute cutoff, and it gives us a concrete reference for a tolerable
cylinder count, namely the DP3 passing range.

The cylinder is a good stand-in for the real filter. On DP3, where the true antibody clash is
known, it predicts clash-free designs (\texttt{af3\_n\_clash\_res}=0) with AUC $0.93$
(\texttt{cylinder\_native\_aware} $0.935$, a little better than the plain count $0.924$;
clash-free versus clashing medians of 5 versus 22). So it stands in well for the filter we
lose on DP4, and the native-aware carve helps the ranking, not just the picture
(\texttt{episcaf\_analysis/cylinder\_surrogate\_auc.py}). This correlation is measured on the
\emph{well-predicted} designs (epitope-chunk RMSD $< \SI{2.5}{\angstrom}$): a design AlphaFold3
folds poorly carries an unreliable epitope placement, which would only add noise to a
clash-versus-surrogate comparison, so it is excluded \emph{from this statistic}. Selecting the
best designs per epitope is a different task and applies no such hard threshold --- it ranks on
the composite score directly (Section~\ref{sec:composite}), which is why the per-island
selection (Section~\ref{sec:selection}) gates only loosely and then ranks.

\begin{figure}[h]\centering
\figorbox{passes_overlay.png}{1.0}
\caption{All designs (grey, left axis) and passing designs (color, right colored axis), both
in true counts, per metric and per row; the dashed line marks the median of the passes. The
twin axes keep the full population and the much smaller passing set each at its own scale, so
the shape of the accepted distribution and where it lands are both visible. DP3 passes are
the four-filter set; DP4 passes are the top 3 per epitope by composite score. The
\emph{Cylinder clashes} column is the \emph{plain} count (\texttt{cylinder\_ca\_clashes}), not
the native-aware carve (\texttt{cylinder\_native\_aware}). Regenerated by
\texttt{episcaf\_analysis/viz/plot\_passes\_overlay.py}.}
\label{fig:overlay}
\end{figure}
